\documentclass[conference]{IEEEtran}
\IEEEoverridecommandlockouts

\usepackage{cite}
\usepackage{amsmath,amssymb,amsfonts}
\usepackage{graphicx}
\usepackage{xcolor}

\begin{document}

\title{Uplink PRB Utilization Forecasting in 5G Networks: Data Preprocessing, Clustering, and Feature Engineering}

\author{
\IEEEauthorblockN{Leander Ziehm}
\IEEEauthorblockA{
Deggendorf Institute of Technology\\
Deggendorf, Germany\\
email
}
\and
\IEEEauthorblockN{Ahnaf Irfan}
\IEEEauthorblockA{
Deggendorf Institute of Technology\\
Deggendorf, Germany\\
email
}
}

\maketitle

\begin{abstract}
This work addresses the problem of forecasting uplink PRB utilization in a live 5G radio access network using high-frequency cell-level measurements. The task involves predicting 3 days ahead at 15-minute resolution (288 steps) and producing probabilistic outputs in the form of quantiles (q10, q50, q90). The dataset contains 226 cells over a 60-day period with no missing values. This report focuses on data preprocessing, exploratory clustering of cells, and feature engineering strategies to support accurate multi-step quantile forecasting.
\end{abstract}

\begin{IEEEkeywords}
5G networks, time series forecasting, PRB utilization, quantile regression, clustering, feature engineering
\end{IEEEkeywords}

\section{Introduction}
The objective is to forecast uplink Physical Resource Block (PRB) utilization in a 5G radio access network. The forecasting horizon is 3 days ahead with 15-minute granularity. The output must be probabilistic, consisting of lower (q10), median (q50), and upper (q90) quantiles.

The dataset contains measurements from 226 cells collected over a 60-day period. Each cell exhibits distinct traffic patterns influenced by user behavior, location type, and network load dynamics.

\section{Dataset Description}
The dataset includes the following variables:
\begin{itemize}
\item Cell identifier
\item Timestamp (15-minute intervals)
\item Uplink PRB utilization percentage
\item Uplink throughput volume
\item Active uplink RRC-connected users
\end{itemize}

The dataset contains 1,301,760 records, with full temporal coverage and no reported missing values. Each cell is expected to contain 5,760 time steps.

\section{Data Preprocessing}

\subsection{Data Integrity Checks}
Although no missing values are reported, the following validations are required:
\begin{itemize}
\item Ensure full 15-minute time continuity per cell
\item Verify uniqueness of (cell, timestamp) pairs
\item Detect and remove duplicate records if present
\end{itemize}

\subsection{Scaling and Normalization}
Feature scaling is required due to differing variable ranges:
\begin{itemize}
\item PRB utilization: percentage-based bounded variable
\item Throughput: high dynamic range variable
\item Active users: count-based discrete variable
\end{itemize}

Robust scaling or standardization is recommended. Log transformation may be applied to throughput if skewness is high. Scaling must be fitted only on training data to prevent leakage.

\subsection{Time Index Construction}
Data must be sorted by cell and timestamp. A strict 15-minute grid should be enforced per cell to ensure consistent sequence modeling.

\section{Cell Clustering Exploration}

\subsection{Motivation}
Cells exhibit heterogeneous behavior depending on geographical and operational conditions. Clustering allows identification of similar traffic patterns, enabling shared modeling strategies and improved generalization.

\subsection{Clustering Approaches}

\subsubsection{Feature-Based K-Means}
Each cell is represented using aggregated statistics such as mean load, variance, peak usage, and correlation metrics. K-Means provides a scalable baseline approach.

\subsubsection{Dynamic Time Warping Clustering}
DTW-based methods capture similarity in temporal patterns even when peak activity is shifted in time. This is useful for daily traffic curves but computationally expensive.

\subsection{Recommended Strategy}
A two-stage approach is suggested:
\begin{itemize}
\item Initial clustering using feature-based K-Means
\item Validation using DTW-based similarity on representative sequences
\end{itemize}

\subsection{Cluster Feature Set}
Useful descriptors include:
\begin{itemize}
\item Mean and maximum PRB utilization
\item Standard deviation of load
\item Throughput statistics
\item Correlation between users and load
\item Daily and weekly variability patterns
\end{itemize}

\section{Feature Engineering}

\subsection{Time-Based Features}
Cyclical encoding is used to represent periodic behavior:
\begin{itemize}
\item Hour of day encoding using sine and cosine transforms
\item Day of week encoding using sine and cosine transforms
\end{itemize}

\subsection{Lag Features}
Temporal dependencies are captured using lagged values:
\begin{itemize}
\item Short-term lags (15 min, 30 min, 1 hour)
\item Medium-term lags (2 hours, 1 day)
\end{itemize}

\subsection{Rolling Statistics}
Rolling windows help smooth noise and capture trends:
\begin{itemize}
\item Rolling mean and standard deviation
\item Rolling min and max
\item Windows: 1 hour, 4 hours, 1 day
\end{itemize}

\subsection{Derived Features}
Additional engineered features include:
\begin{itemize}
\item Throughput per active user ratio
\item Rate of change in PRB utilization
\item Interaction terms between traffic variables
\end{itemize}

\section{Modeling Implications}
A global forecasting model across all cells is recommended, optionally enhanced with cluster identifiers or embeddings. The model should support multi-step forecasting over 288 future time steps and output quantiles using a quantile loss function.

\section{Workflow Summary}
The recommended pipeline is:
\begin{itemize}
\item Data validation and cleaning
\item Feature scaling and preprocessing
\item Feature engineering (lags, rolling, time encoding)
\item Clustering analysis of cells
\item Model training using shared representation
\item Evaluation against baseline per-cell models
\end{itemize}

\section{Conclusion}
The dataset provides a strong foundation for time series forecasting with full temporal coverage and multiple informative features. Effective preprocessing, feature engineering, and clustering are key to building a scalable and accurate forecasting system for 5G network load prediction.

TODO.
\end{document}